% Part: first-order-logic
% Chapter: sequent-calculus
% Section: quantifier-rules

\documentclass[../../../include/open-logic-section]{subfiles}

\begin{document}

\olfileid{fol}{seq}{qrl}

\olsection{Aturan Kuantor}

\subsection{Aturan kanggo $\lforall$}

\begin{defish}
\Axiom$ !A(t), \Gamma \fCenter \Delta$
\RightLabel{\LeftR{\lforall}}
\UnaryInf$ \lforall[x][!A(x)],\Gamma \fCenter \Delta$
\DisplayProof
\hfill
\Axiom$ \Gamma \fCenter \Delta, !A(a) $
\RightLabel{\RightR{\lforall}}
\UnaryInf$ \Gamma \fCenter \Delta, \lforall[x][!A(x)]$
\DisplayProof
\end{defish}

Ing \LeftR{\lforall}, $t$ iku term tertutup (yaiku term tanpa
variabel). Ing \RightR{\lforall}, $a$~iku !!a{constant} kang ora kena
muncul ing ngendi wae ing sekuen ngisor saka aturan
\RightR{\lforall}. Kita ngarani $a$ \emph{eigenvariabel} saka inferensi
\RightR{\forall}.\footnote{Kita nggunakake istilah “eigenvariabel”
sanajan $a$ ing aturan ndhuwur iku !!a{constant}. Iki amarga alasan
sejarah.}

\subsection{Aturan kanggo $\lexists$}

\begin{defish}
\Axiom$ !A(a), \Gamma \fCenter \Delta $
\RightLabel{\LeftR{\lexists}}
\UnaryInf$ \lexists[x][!A(x)], \Gamma \fCenter \Delta$
\DisplayProof
\hfill
\Axiom$ \Gamma \fCenter \Delta, !A(t) $
\RightLabel{\RightR{\lexists}}
\UnaryInf$ \Gamma \fCenter \Delta, \lexists[x][!A(x)]$
\DisplayProof
\end{defish}

Maneh, $t$~iku term tertutup, lan $a$~iku !!a{constant} kang ora
muncul ing sekuen ngisor saka aturan \LeftR{\lexists}. Kita ngarani
$a$ \emph{eigenvariabel} saka inferensi \LeftR{\lexists}.

Syarat supaya eigenvariabel ora muncul ing sekuen ngisor saka
inferensi \RightR{\lforall} utawa \LeftR{\lexists} diarani
\emph{syarat eigenvariabel}.

\begin{explain}
Elinga konvensi yen nalika $!A$ iku !!a{formula} kanthi
!!{variable}~$x$ bebas, kita nuduhake iki kanthi nulis~$!A(x)$. Ing
konteks kang padha, $!A(t)$ iku cekakan saka~$\Subst{!A}{t}{x}$. Mula
kita uga bisa nulis aturan \RightR{\lexists} kaya mangkene:
\begin{prooftree}
  \Axiom$\Gamma \fCenter \Delta, \Subst{!A}{t}{x}$
  \RightLabel{\RightR{\lexists}}
  \UnaryInf$\Gamma \fCenter \Delta, \lexists[x][!A]$
\end{prooftree}
Gatekna yen $t$ bisa wae wis muncul ing~$!A$, upamane $!A$~bisa
awujud~$\Atom{\Obj P}{t,x}$. Dadi, nginferensi $\Gamma \Sequent \Delta,
\lexists[x][\Atom{\Obj P}{t,x}]$ saka~$\Gamma \Sequent \Delta,
\Atom{\Obj P}{t,t}$ iku panrapan \RightR{\lexists} kang bener---kita
bisa “ngganti” siji utawa luwih, lan ora kudu kabeh, muncule~$t$
ing premis nganggo !!{variable} terikat~$x$. Nanging, syarat
eigenvariabel ing \RightR{\lforall} lan~\LeftR{\lexists} nuntut supaya
!!{constant}~$a$ ora muncul ing~$!A$. Mula, kita ora bisa kanthi bener
nginferensi $\Gamma \Sequent \Delta,
\lforall[x][\Atom{\Obj P}{a,x}]$ saka $\Gamma \Sequent \Delta,
\Atom{\Obj P}{a,a}$ nganggo~$\RightR{\lforall}$.
\end{explain}

\begin{explain}
Ing \RightR{\lexists} lan \LeftR{\lforall} ora ana watesan tumrap
term~$t$. Kosok baline, ing aturan \LeftR{\lexists} lan
\RightR{\lforall}, syarat eigenvariabel nuntut supaya
!!{constant}~$a$ ora muncul ing ngendi wae ing njabane~$!A(a)$ ing
sekuen ndhuwur. Syarat iki perlu kanggo njamin sistem tetep sah, yaiku
mung !!{derive}s sekuen kang valid. Tanpa syarat iki, inferensi iki
bakal diidini:
\begin{prooftree}
  \Axiom$!A(a) \fCenter !A(a)$
  \RightLabel{*\LeftR{\lexists}}
  \UnaryInf$\lexists[x][!A(x)] \fCenter !A(a)$
  \RightLabel{\RightR{\lforall}}
  \UnaryInf$\lexists[x][!A(x)] \fCenter \lforall[x][!A(x)]$
  \DisplayProof\bottomAlignProof
  \qquad
  \Axiom$!A(a) \fCenter !A(a)$
  \RightLabel{*\RightR{\lforall}}
  \UnaryInf$!A(a) \fCenter \lforall[x][!A(x)]$
  \RightLabel{\LeftR{\lexists}}
  \UnaryInf$\lexists[x][!A(x)] \fCenter \lforall[x][!A(x)]$
\end{prooftree}
Nanging, $\lexists[x][!A(x)] \Sequent \lforall[x][!A(x)]$ ora valid.
\end{explain}

\end{document}
